\documentclass{article}
\usepackage[T1]{fontenc}
\usepackage{lmodern}
\usepackage{graphicx}
\usepackage{amsmath,amssymb}
\usepackage[a4paper,margin=2cm]{geometry}
\usepackage[absolute,overlay]{textpos}
\setlength{\TPHorizModule}{1mm}
\setlength{\TPVertModule}{1mm}
\usepackage[indonesian]{babel}
\usepackage{hyperref}
\setlength{\emergencystretch}{2em}
% unit: Halaman 28

\begin{document}

% === Halaman 28 (llm) ===

\begin{itemize}
    \item Bagaimana jika terdapat dua atau lebih nilai $k$ yang menghasilkan pesan-pesan bermakna?
\end{itemize}

Contoh: Misalkan kriptogram \texttt{HSPPW} menghasilkan dua kemungkinan kunci yang potensial, yaitu:

\begin{align*}
    k = 4 & \text{ menghasilkan pesan } \texttt{dolls} \quad (\text{boneka}) \\
    k = 11 & \text{ menghasilkan } \texttt{wheel} \quad (\text{roda})
\end{align*}

Nilai $k$ mana yang benar?

Jika kasusnya demikian, maka lakukan dekripsi terhadap potongan cipherteks lain tetapi cukup menggunakan $k = 4$ dan $k = 11$ agar dapat disimpulkan kunci mana yang benar.

\end{document}
